\documentclass[french]{article}

%Chargement des paquets
\usepackage{amsmath}
\usepackage{amsfonts}
\usepackage{amssymb}
\usepackage{enumerate}
\usepackage{mathtools}
\usepackage{bbm}
\usepackage{xparse, etoolbox}
\usepackage{enumerate}
\usepackage{enumitem}
\usepackage{mathabx}
\usepackage{babel}[french]

%environment
\usepackage{amsthm}
\usepackage{keytheorems}
\usepackage{hyperref} 

%Interface théorème
\newkeytheorem{lem}[
	name = Lemme
]

\newkeytheorem{prop}[
	name = Proposition
]

\newkeytheorem{thm}[
	name = Théorème
]

\newkeytheorem{coro}[
	name = Corollaire
]

\renewcommand*{\proofname}{Démonstration}

\theoremstyle{definition}
\newtheorem{defn}{Définition}

\theoremstyle{definition}
\newtheorem*{nota}{Notation}

\theoremstyle{definition}
\newtheorem*{limi}{Liminaire}

\theoremstyle{remark}
\newtheorem{rmq}{Remarque}

\theoremstyle{plain}
\newtheorem*{fait}{Fait}


%Ensembles de nombres
\newcommand{\N} {\mathbb{N}}
\newcommand{\Ne}{\N^\ast}
\newcommand{\Z} {\mathbb{Z}}
\newcommand{\Q} {\mathbb{Q}}
\newcommand{\R} {\mathbb{R}}
\newcommand{\Rb}{\overline{\mathbb{R}}}
\newcommand{\Rp}{\R_+}
\newcommand{\Rm}{\R_-}
\newcommand{\K} {\mathbb{K}}
\newcommand{\Ket} {\mathbb{K^{\ast}}}
\newcommand{\Cx}{\mathbb{C}}

%Ensembles, tribus
\newcommand{\A}{\mathbb{A}}
\renewcommand{\P}{\mathcal{P}}
\newcommand{\T}{\mathcal{T}}
\newcommand{\F}{\mathcal{F}}
\newcommand{\C} {\mathcal{C}}
\newcommand{\ouv}{\mathcal{O}}
\newcommand{\B}{\mathcal{B}}
\newcommand{\M}{\mathcal{M}}
\newcommand{\etag}{\mathcal{E}}
\newcommand{\etagOp}{\etag^{+}(\Omega)}
\newcommand{\etagO}{\etag(\Omega)}
%%Lp avant quotientage
\newcommand{\Lic}{\mathcal{L}^1}
\newcommand{\Lpc}{\mathcal{L}^p}
\newcommand{\Linfc}{\mathcal{L}^{\infty}}
\newcommand{\Lifull}{\Lic(\Omega, \B, m)}
\newcommand{\Lpfull}{\Lpc(\Omega, \B, m)}
\newcommand{\Linffull}{\Linfc(\Omega, \B, m)}
%%Lp après quotientage
\newcommand{\Li}{L^1}
\newcommand{\Ld}{L^2}
\newcommand{\Lp}{L^p}
\newcommand{\Linf}{L^{\infty}}
\newcommand{\Listd}{\Li(\Omega, m)} 
\newcommand{\Ldstd}{\Ld(\Omega, m)} 
\newcommand{\Lpstd}{\Lp(\Omega, m)} 

%Quelques opérateurs
\DeclareMathOperator{\codim}{codim}
\DeclareMathOperator{\rg}{rg}
\DeclareMathOperator{\tr}{tr}
\DeclareMathOperator{\vect}{Vect}
\DeclareMathOperator{\ima}{Im}
\DeclareMathOperator{\id}{Id}
\DeclareMathOperator{\sign}{sign}
\DeclareMathOperator{\ext}{ext}
\newcommand{\ind}{\mathbbm{1}}
\newcommand*{\spr}[2]{\left(\,#1\,\middle|\,#2\,\right)}
\newcommand{\interior}[1]{%
  {\kern0pt#1}^{\mathrm{o}}%
}
\DeclareMathOperator{\card}{card}
\DeclareMathOperator{\ppcm}{ppcm}

%Commandes de pure flemme
\newcommand{\veps}{\varepsilon}
\newcommand{\intab}{\int_{a}^{b}}
\newcommand{\intR}{\int_{-\infty}^{\infty}}
\newcommand{\intinf}{\int_{- \infty}^{+ \infty}}
\newcommand{\equi}{\Leftrightarrow}
\newcommand{\Kev}{$\K$-espace vectoriel }
\newcommand{\Kevs}{$\K$-espaces vectoriels }
\newcommand{\Rev}{$\R$-espace vectoriel }
\newcommand{\Revs}{$\R$-espaces vectoriels }
\newcommand{\Cev}{$\Cx$-espace vectoriel }
\newcommand{\Cevs}{$\Cx$-espaces vectoriels }
\newcommand{\evn}{(E, \norm{.})}
\newcommand{\interf}[2]{[#1,#2]}
\newcommand{\limb}{\lim\limits_}
\newcommand{\lebn}{\lambda_n}
\newcommand{\pp}{\text{~presque partout}}
\newcommand{\derivp}[2]{\frac{\partial #1}{\partial #2}}
\newcommand{\famiu}{(u_i)_{i \in I}}
\newcommand{\sev}{\text{~s.e.v.}}
\newcommand{\Mnp}{M_{n,p}(\K)}
\newcommand{\Mn}{M_{n}(\K)}
\newcommand{\tq}{\text{~t.q.~}}
\newcommand{\mq}{~montrer~que~}
\newcommand{\et}{\quad \text{et} \quad}
\newcommand{\ou}{\text{~ou~}}
\newcommand{\Kx}{\K[X]}


%Commandes de gars chelou
\renewcommand{\subset}{\subseteq}

%Commande restriction, ty duckduckgo
\def\restr#1#2{\mathchoice
              {\setbox1\hbox{${\displaystyle #1}_{\scriptstyle #2}$}
              \restraux{#1}{#2}}
              {\setbox1\hbox{${\textstyle #1}_{\scriptstyle #2}$}
              \restraux{#1}{#2}}
              {\setbox1\hbox{${\scriptstyle #1}_{\scriptscriptstyle #2}$}
              \restraux{#1}{#2}}
              {\setbox1\hbox{${\scriptscriptstyle #1}_{\scriptscriptstyle #2}$}
              \restraux{#1}{#2}}}
\def\restraux#1#2{{#1\,\smash{\vrule height .8\ht1 depth .85\dp1}}_{\,#2}} 

%Quelques normes
\DeclarePairedDelimiter\abs{\lvert}{\rvert}
\DeclarePairedDelimiter\labs{\bigl\lvert}{\bigr\rvert}
\DeclarePairedDelimiter\norm{\lVert}{\rVert}
\DeclarePairedDelimiterXPP\normi[1]{}\lVert\rVert{_1}{\ifblank{#1}{\:\cdot\:}{#1}}
\DeclarePairedDelimiterXPP\normd[1]{}\lVert\rVert{_2}{\ifblank{#1}{\:\cdot\:}{#1}}
\DeclarePairedDelimiterXPP\normp[1]{}\lVert\rVert{_p}{\ifblank{#1}{\:\cdot\:}{#1}}
\DeclarePairedDelimiterXPP\normq[1]{}\lVert\rVert{_q}{\ifblank{#1}{\:\cdot\:}{#1}}
\DeclarePairedDelimiterXPP\norminf[1]{}\lVert\rVert{_\infty}{\ifblank{#1}{\:\cdot\:}{#1}}

%Bracket en angle
\DeclarePairedDelimiter\angl{\langle}{\rangle}

%Set, parenthèses
\newcommand{\set}[1]{\{ #1 \}}
\newcommand{\bigset}[1]{\bigl\{ #1 \bigr\}}
\newcommand{\Bigset}[1]{\Bigl\{ #1 \Bigr\}}
\newcommand{\bigpar}[1]{\bigl( #1 \bigr)}
\newcommand{\Bigpar}[1]{\Bigl( #1 \Bigr)}

%Opitmisation des opérateurs de comparaison
\DeclarePairedDelimiter\tio{o (}{)}
\DeclarePairedDelimiter\gro{O (}{)}


\begin{document}

Dans tout ce document, les suites auront pour support $\N$ tout entier. Il est clair que les résultats restent valables pour tout support 
$S \subset \N$ infini.

\section{Définition et premiers résultats}

\begin{defn}[Série]
Soit $(a_n)_\N \subset \K$ une suite numérique avec $\K = \R$ ou $\Cx$. 
On définit la suite $(S_n)_\N$ par :
\[ \forall N \in \N, S_N = \sum_{n=0}^N a_n . \]
On appelle $(S_n)_\N$ la série de terme général $a_n$ et on la note $\sum a_n$.
\end{defn}

\begin{defn}[Convergence]
Soit $(a_n)_\N \subset \K$ une suite numérique avec $\K = \R$ ou $\C$. 
On dit que la série $\sum a_n$ converge si, et seulement si, la suite $(S_n)_\N$ converge.
Dans ce cas, la limite de $(S_n)_\N$ est notée $\sum_{n=0}^{\infty} a_n$.
\end{defn}

\begin{thm}
Si la série $\sum a_n$ converge alors la suite $(a_n)_\N$ converge vers 0.
\end{thm}

\begin{proof}
La suite $(S_n)_\N$ converge, elle est donc de Cauchy, ce qui implique que $(a_n)_\N$ est de limite nulle.
\end{proof}

\begin{coro}
La série harmonique $\sum \frac1{n}$ est divergente.
\end{coro}

\begin{proof}
Soit $S_N = \sum_{n=1}^N\frac{1}{n} $, on a 
\[ S_{2N} - S_N = \sum_{n=N+1}^{2N} \dfrac{1}{n} \geq N \dfrac1{2N} = \dfrac12 , \]
donc la suite $(S_N)_{\N \geq 1}$ n'est pas de Cauchy, elle n'est pas convergente.
\end{proof}

\section{Propriétés fondamentales}

\begin{thm}[Linéarité]
L'ensemble des séries convergentes est un \Kev.
\end{thm}

\begin{proof}
Cela revient à dire que l'ensemble des suites convergentes est un \Kev, ce qui est vrai.
\end{proof}

\begin{coro}
Soit $\sum a_n$ une série convergente et $\sum v_n$ une série divergente, alors $\sum (a_n+v_n)$ diverge.
\end{coro}

\begin{proof}
Sinon, la série $\sum( a_n + v_n - a_n) = \sum v_n$ serait convergente.
\end{proof}

\begin{thm}[Critères de comparaison]
Soient $(a_n)_\N$ et $(b_n)_\N$ deux suites réelles positives.
\begin{enumerate}[label=(\roman*)]
\item s'il existe $n_0 \in \N$ tel que $\forall n \geq n_0, a_n \leq b_n$ et si $\sum b_n$ converge alors $\sum a_n$ converge ;
\item s'il existe $n_0 \in \N$ tel que $\forall n \geq n_0, a_n \geq b_n$ et si $\sum b_n$ diverge alors $\sum a_n$ diverge.
\end{enumerate}
\end{thm}

\begin{proof}
On note $S_N = \sum_{n=0}^N a_n$ et $T_N = \sum_{n=0}^N a_n$. 
Pour alléger la démonstration, on suppose que $n_0 = 0$ (le cas $n_0$ quelconque se traite de la même manière).
\begin{enumerate}[label=(\roman*)]
\item On remarque que les suite $(S_n)_\N$ et $(T_N)_\N$ sont croissantes, on peut donc écrire :
\[ \forall n \in \N, S_n \leq \sum_{n=0}^{\infty} b_n, \] 
donc $(S_n)_\N$ est croissante et majorée, elle est donc convergente.
\item La proposition $\forall n \in \N, b_n \leq a_n$, $\sum b_n$ diverge et $\sum a_n$ converge est fausse selon (i), d'où le résultat.
\end{enumerate}
\end{proof}

\begin{thm}[Séries équivalentes]
Soient $(a_n)_\N$ et $(b_n)_\N$ deux suites réelles. Si $a_n \sim_{\infty} b_n$ et si $a_n$ est de signe constant à partir d'un certains rang,
alors $\sum a_n$ et $\sum b_n$ sont de même nature.
\end{thm}

\begin{proof}
Suppons qu'il existe $n_0 \in \N$ tel que $\forall n \geq n_0, a_n \geq 0$ (la démonstration est la même pour $a_n$ négatif à partir d'un
certains rang). Notons $(u_n)_\N$ la suite réelle telle que :
\[ a_n = u_n b_n \et \lim u_n =1 ,\]
il existe $n_1 \in \N$ tel que $\forall n \geq n_1, \frac12 \leq u_n \leq \frac32$. On note $N = \max(n_0, n_1)$, pour $n \geq N$ on a 
$b_n \geq 0$ (car $a_n = u_n b_n, a_n \geq 0$ et $u_n \geq 0$) et donc :
\[ \forall n \geq N, \frac12 b_n \leq a_n \leq \frac32 b_n , \]
ce qui permet de conclure.
\end{proof}

\begin{defn}[Convergence absolue]
Soit $(a_n)_\N \subset \K$ une suite numérique. On dit que $\sum a_n$ converge absoluement si la série $\sum \abs{a_n}$ converge. 
\end{defn}

\begin{thm}
Soit $(a_n)_\N \subset \K$ une suite numérique. Si la série $\sum a_n$ converge absoluement alors elle converge.
\end{thm}

\begin{proof}
Supposons tout d'abord que $(a_n)_\N$ est une suite à valeurs réelles. On note :
\[ a_n^{+} = \left \{ \begin{array}{c c} a_n & \text{ si } a_n \geq 0 \\ 0 & \text{ sinon} \end{array} \right.
\et a_n^{-} = \left \{ \begin{array}{c c} 0 & \text{ si } a_n \geq 0 \\ -a_n & \text{ sinon} \end{array} \right. \]
Les séries $\sum a_n^{+}$ et $\sum a_n^{-}$ sont à termes positifs et :
\[ \forall n \in \N, a_n^{+} \leq \abs{a_n} \et a_n^{-} \leq \abs{a_n} , \]
elles sont donc convergentes. Or $\sum a_n = \sum a_n^{+}- \sum a_n^{-}$ d'où le résultat dans le cas réel. \newline

Dans le cas complexe, on procède de la même façon en écrivant successivement :
\[ \abs{\Re(z)} \leq \abs{z} \et \abs{\Im(z)} \leq \abs{z} ; \]
on applique ensuite le cas réel aux deux parties ce qui permet de conclure.
\end{proof}

\begin{thm}[Généralisation]
Soit $(E, \norm{.})$ un espace vectoriel normé complet et soit $(a_n)_\N$ une suite d'éléments de $E$.
Si $\sum \norm{a_n}$ converge alors $\sum a_n$ converge.
\end{thm}

\begin{proof}
Comme $\sum \norm{a_n}$ converge, la suite $\sum_{n=0}^{N} \norm{a_n}$ vérifie le critère de Cauchy. 
Soit $\veps >0$, il existe $n_0 \in \N$ tel que :
\[ \forall N \geq n_0, \forall k \geq 0, \sum_{n=N+1}^{N+k} \norm{a_n} < \veps . \]
Donc $\norm*{\sum_{n=N+1}^{N+k} {a_n}} \leq \sum_{n=N+1}^{N+k} \norm{a_n} < \veps$, ainsi la suite des sommes partielles 
est de Cauchy, elle est donc convergente par complétude de $E$.
\end{proof}

\section{Comparaison série-intégrale}

Soit $f : \Rp \to \R$ une fonction continue, positive et décroissante. Le but de cette section est d'étudier les séries de la forme :
$\sum f(n)$.

\begin{thm}
Soit $f : \Rp \to \R$ une fonction continue, positive et décroissante, la série $\sum f(n)$ converge si, et seulement si, la suite 
$\Bigpar{\int_{a}^N f(x) dx}_\N$ admet une limite lorsque $N \to \infty$ (avec $a \in \Rp$ quelconque).
\end{thm}

\begin{proof}
Soit $n \in \N$, comme $f$ est décroissante et positive, on peut écrire ;
\[ \forall x \in [n ; n+1], f(n+1) \leq f(x) \et \forall x \in [n ; n+1] f(x) \leq f(n) , \]
d'où il résulte que :
\[ \int_1^{N+1} f(x) dx \leq \sum_{n=1}^N f(n) \leq \int_0^{N} f(x) dx . \]
Ce qui permet d'obtenir le résultat.
\end{proof}

\begin{coro}[Séries de Riemann]
Soit $\alpha \in \R \setminus \{1\}$, on appelle série de Riemann la série $\sum_{n \geq 1} \frac{1}{n^{\alpha}}$.
Cette série converge si, et seulement si, $\alpha > 1$.
\end{coro}

\begin{proof}
Si $\alpha \leq 0$, la suite $\bigpar{\frac{1}{n^{\alpha}}}_\N$ ne tend pas vers 0 donc la série diverge.
Sinon, la fonction $x \mapsto \frac{1}{x^{\alpha}}$ est continue, positive et décroissante sur $]0, \infty[$. De plus :
\[ \int_1^N \dfrac{1}{n^{\alpha}} = \left [ \dfrac{x^{1-\alpha}}{1-\alpha} \right ]_1^N = \dfrac{N^{1-\alpha}}{1-\alpha} - \dfrac{1}{1-\alpha} . \]
Ainsi, l'intégrale (et donc la série) converge si, et seulement si $\alpha > 1$.
\end{proof}

\begin{coro}[Série harmonique]
Pour $N \geq 1$ on a :
\[ \sum_{n=1}^N \dfrac1{n} \sim_{\infty} \ln{N} . \]
\end{coro}

\begin{proof}
On reprend l'encadrement de la démonstration du théorème, on a :
\[ \int_2^{N+1} \dfrac1{x} dx \leq \sum_{n=2}^N \dfrac1{n} \leq \int_1^{N} \dfrac1{x} dx . \]
soit, en ajoutant 1 à chaque terme de l'inégalité :
\[ 1 + \ln(N+1) - \ln{2} \leq \sum_{n=1}^N \dfrac1{n} \leq 1 + \ln{N} , \]
d'où on tire l'équivalence.
\end{proof}

\section{Quelques critères de convergence}

Dans toute section, on suppose que $(a_n)_\N$ est une suite réelle positive. 
On pourra étendre les critères aux suites à valeurs dans $\K$ en utilisant l'implication convergence absolue $\Rightarrow$ convergence.

\begin{thm}[Critère de D'Alembert]
On suppose que $\lim \dfrac{a_{n+1}}{a_n} = L$. Alors :
\begin{enumerate}[label=(\roman*)]
\item si $L < 1$, la série $\sum a_n$ converge ;
\item si $L>1$, la série $\sum a_n$ diverge ;
\item si $L=1$, on ne peut pas conclure.
\end{enumerate}
\end{thm}

\begin{proof}
\begin{enumerate}[label=(\roman*)]
\item supposons que $L<1$, alors il existe $q < 1$ et $N \in \N$ tels que :
\[ \forall n \geq N, \dfrac{a_{n+1}}{a_n} < q . \]
On peut donc écrire, pour tout $k \in \N, a_{N+k} < a_N q^k$ et on remarque que :
\[ \sum_{k=0}^n a_{N+k} < a_N \sum_{k=0}^n q^k , \]
ce qui permet de conclure par comparaison.
\item si $L>1$, par un raisonnement similaire on montre que pour tout $k \in \N, a_{N+k} > a_N$. 
La série est donc trivialement divergente.
\item Prenons la suite définie sur $\Ne$ par $a_n = \frac1{n^{\alpha}}$, on a $a_{n+1}/a_n \to 1$ quelle que soit la valeur de $\alpha$
et pourtant $\sum a_n$ diverge si $\alpha \leq 1$. 
\end{enumerate}
\end{proof}

\begin{rmq}
On peut énoncer une version alternative de ce critère avec des hypothèses légèrement plus faibles, à savoir :
\begin{enumerate}[label=(\roman*)]
\item si $\dfrac{a_{n+1}}{a_n} < q < 1$ à partir d'un certains rang, alors $\sum a_n$ converge ;
\item si $\dfrac{a_{n+1}}{a_n} > 1$ à partir d'un certains rang, alors $\sum a_n$ diverge.
\end{enumerate}
\end{rmq}

\begin{thm}[Critère de Cauchy]
On suppose que $\lim (a_n)^{1/n} = L$. Alors :
\begin{enumerate}[label=(\roman*)]
\item si $L < 1$, la série $\sum a_n$ converge ;
\item si $L>1$, la série $\sum a_n$ diverge ;
\item si $L=1$, on ne peut pas conclure.
\end{enumerate}
\end{thm}

\begin{proof}
La preuve est identique à celle du critère de D'Alembert.
\end{proof}

\begin{rmq}
Comme précédemment, on peut énoncer une version alternative de ce critère avec des hypothèses légèrement plus faibles, à savoir :
\begin{enumerate}[label=(\roman*)]
\item si $(a_n)^{1/n} < q < 1$ à partir d'un certains rang, alors $\sum a_n$ converge ;
\item si $(a_n)^{1/n} > 1$ à partir d'un certains rang, alors $\sum a_n$ diverge.
\end{enumerate}
\end{rmq}

\begin{thm}[Critère de Riemann]
\begin{enumerate}[label=(\roman*)]
\item si $\lim n^{\alpha} a_n = 0$ pour $\alpha > 1$ alors $\sum a_n$ converge ;
\item si $\lim n^{\alpha} a_n = \infty$ pour $\alpha < 1$ alors $\sum a_n$ diverge.
\end{enumerate}
\end{thm}

\begin{proof}
La structure de la preuve reste la même.
\begin{enumerate}[label=(\roman*)]
\item si $\lim n^{\alpha} a_n = 0$ alors il existe un rang $N \in \N$ à partir duquel $\forall n \geq N, n^{\alpha} a_n \leq 1$ 
et on obtient le résultat par comparaison avec une série de Riemann.
\item On démontre de même.
\end{enumerate}
\end{proof}

\begin{rmq}
La version avec les hypothèses légèrement plus faibles est :
\begin{enumerate}[label=(\roman*)]
\item si $n^{\alpha} a_n \leq 1$ à partir d'un certains rang pour $\alpha > 1$, alors $\sum a_n$ converge ;
\item si $n^{\alpha} a_n \geq 1$ à partir d'un certains rang pour $\alpha < 1$, alors $\sum a_n$ diverge.
\end{enumerate}
\end{rmq}

\begin{thm}[Séries de Bertrand]
Soit $\alpha, \beta \in \R$, on appelle série de Bertrand la série 
\[ \sum_{n \geq 2} \dfrac{1}{n^{\alpha}(\ln{n})^{\beta}} . \]
Cette série converge si, et seulement si, ($\alpha > 1$) ou ($\alpha = 1$ et $\beta > 1$).
\end{thm}

\begin{proof}
Pour $\alpha > 1$, on peut trouver un réel $\gamma$ tel que $1 < \gamma < \alpha$ et :
\[ \lim_{n \to \infty} n^{\gamma} \dfrac{1}{n^{\alpha}(\ln{n})^{\beta}} = 
\lim_{n \to \infty} \dfrac{1}{n^{\alpha-\gamma}(\ln{n})^{\beta}} = 0 \]

Pour $\alpha <1$, on peut trouver un réel $\gamma$ tel que $\alpha < \gamma < 1$ et :
\[ \lim_{n \to \infty} n^{\gamma} \dfrac{1}{n^{\alpha}(\ln{n})^{\beta}} = 
\lim_{n \to \infty} \dfrac{n^{\alpha-\gamma}}{(\ln{n})^{\beta}} = \infty \]

Pour $\alpha = 1$ on considère plusieurs cas.
Si $\beta \leq 0$, on a $\forall n \geq 3, \frac{1}{n(\ln{n})^{\beta}} > \frac{1}{n}$ donc la série diverge.
Sinon, la fonction $x \mapsto \dfrac{1}{x(\ln{x})^\beta}$ est continue, positive et décroissante sur $]1, \infty[$, et
\[ \arraycolsep=1.4pt\def\arraystretch{2.2}
\int_2^N \dfrac{1}{x(\ln{x})^{\beta}} = \left \{ \begin{array}{c c}
\left [ \dfrac{(\ln{x})^{1-\beta}}{1-\beta} \right ]_2^N  & \text{ si } \beta \neq 0 \\
\left [ \ln(\ln(x)) \right ]_2^N & \text{ sinon}
\end{array} \right.,  \]
qui converge si, et seulement si, $\alpha >1$.
\end{proof}

\section{Séries alternées}

Pour une série $\sum a_n$ convergente vers $S \in \R$, on note, $R_N = S - S_n$.

\begin{defn}[Série alternée]
Soit $(a_n)_\N$ une suite réelle. On dit que la série $\sum a_n$ est alternée si pour tout $n \in \N, a_{n+1}a_n < 0$,
c'est-à-dire si pour tout $n \in \N, a_n = (-1)^n u_n$ avec $(u_n)$ de signe constant. 
\end{defn}

\begin{thm}[Règle de Leibniz]
Soit $\sum a_n$ une série alternée. Si la suite $(\abs{a_n})_\N$ est décroissante et tend vers 0 alors $\sum a_n$ converge.
\end{thm}

\begin{proof}
Quitte à changer le signe de la série, on suppose que $a_0 > 0$.
On a : $0 < a_0 + a_1 < a_0$ puis $a_0+a_1 < a_0 + a_1 + a_2$ etc.
Notons, pour $n \in \N, p_n = S_{2n}$ et $i_n = S_{2n+1}$ et montrons que les suites $(p_n)_\N$ et $(i_n)_\N$ sont adjacentes.
Il est clair que $(p_n)_\N$ est décroissante : $p_{n+1} - p_n = a_{2n+2} + a_{2n+1} < 0$
et de même $(i_n)_\N$ est croissante. De plus $p_n -  i_n = - a_{2n+1}$ qui tend vers 0.
Les suites $p_n$ et $i_n$ sont donc adjacentes, on en déduit qu'elles convergent vers la même limite $S$ et, par suite, $S_n \to S$.
\end{proof}

\begin{thm}[Majoration du reste]
Soit $\sum a_n$ une série alternée avec $(\abs{a_n})_\N$ décroissante et qui tend vers 0.
On a alors :
\[ \forall N \in \N, \abs*{R_N} = \abs*{\sum_{n=N+1}^\infty a_n} \leq \abs{a_{N+1}} . \]
\end{thm}

\begin{proof}
Soit $n \in \N$, on va montrer que, pour tout $p \in \N, a_n + a_{n+1} + \dots + a_{n+p}$ est du même signe que $a_n$.
En effet, tous les couples de la forme $a_{n+2k} + a_{n+2k+1}$ pour $k \geq 0$ sont du même signe que $a_n$. \\
Ainsi, dans la somme $a_n + a_{n+1} + \dots + a_{n+p}$ on regroupe tous les termes par paires, chaque paire est du même signe que 
$a_n$. De plus, si le dernier terme est isolé, son rang est nécessairement de la forme $a_{n+2k}$ avec $k \geq 0$ 
et son signe est le même que celui de $a_n$. 
Par passage à la limite, on conclut que pour tout $N \in \N, R_N$ est du même signe que $a_{N+1}$.
Quitte à changer le signe de la série, on peut supposer que $a_{N+1} > 0$. On a alors $R_N > 0$ et $R_{N+1} < 0$
or $R_N = a_{N+1} + R_{N+1}$ d'où $R_N \leq a_{N+1}$.
\end{proof}

\begin{lem}[Transformation d'Abel]
Pour toute suite $u=(u_n)_\N$ réelle ou complexe, on définit une suite $Du$ par :
\[ \forall n \in \N, Du_n = u_{n+1} - u_n . \]
Pour deux suites réelles ou complexes $(u_n)_\N$ et $(v_n)_\N$ on a :
\[ \sum_{k=0}^{n-1} Du_k v_{k+1} = (u_n v_n - u_0 v_0) - \sum_{k=0}^{n-1} u_k Dv_k . \] 
\end{lem}

\begin{proof}
Il suffit de faire le calcul :
\begin{align*}
\sum_{k=0}^{n-1} Du_k v_{k+1} &= \sum_{k=0}^{n-1} u_{k+1} v_{k+1} - \sum_{k=0}^{n-1} u_k v_{k+1} \\
& = \sum_{k=1}^{n} u_{k} v_{k} - \sum_{k=0}^{n-1} u_k v_{k+1} \\
& = (u_n v_n - u_0 v_0) - (\sum_{k=0}^{n-1} u_k v_{k+1} - \sum_{k=0}^{n-1} u_k v_{k} \\
& = (u_n v_n - u_0 v_0) - \sum_{k=0}^{n-1} u_k Dv_k
\end{align*}
\end{proof}

\begin{thm}[Théorème d'Abel]
Soit $(a_n)_\N$ une suite réelle décroissante qui tend vers 0 (donc positive).
Soit $(b_n)_\N$ une suite réelle ou complexe. On note $B_N = \sum_{n=0}^N b_n$ et on suppose que la suite $(B_n)_\N$ est bornée.
Alors la série $\sum a_n b_n$ converge.
\end{thm}

\begin{proof}
On remarque que $\forall N \in \N, DB_k = b_{k+1}$. On a donc :
\begin{align*}
\sum_{k=0}^n a_k b_k & = a_0 b_0 + \sum_{k=0}^{n-1} a_{k+1} b_{k+1} \\
& =  a_0 b_0 + \sum_{k=0}^{n-1} a_{k+1} DB_k \\
& = a_0 b_0 + B_n a_n - B_0 a_0 - \sum_{k=0}^{n-1} Da_k B_k \text{ transformation d'Abel} \\
& = B_n a_n - \sum_{k=0}^{n-1} (a_{k+1} - a_k) B_k
\end{align*}
Comme $(B_n)_\N$ est bornée, on peut écrire $\forall n \in \N, \abs{B_n} \leq B$. On a :
\begin{align*}
\abs*{ \sum_{k=0}^{n-1} (a_{k+1} - a_k) B_k} & \leq B \sum_{k=0}^{n-1} (a_k - a_{k+1}) \\
& \leq B (a_0 - a_n) 
\end{align*}
Donc la série $\sum (a_{k+1} - a_k) B_k$ est absolument convergente. De plus, la suite $(B_n a_n)_\N$ converge vers 0. 
On peut donc conclure que la série $\sum a_n b_n$ converge.
\end{proof}

\section{Produit de Cauchy}

\begin{defn}
Soient $\sum a_n$ et $\sum b_n$ deux séries à valeurs dans $\Cx$. On pose :
\[ \forall n \in \N, c_n = \sum_{k=0}^n a_k b_{n-k} . \]
La série $\sum c_n$ est appelée produit de Cauchy des séries $\sum a_n$ et $\sum b_n$.
\end{defn}

\begin{thm}
Si $\sum a_n$ et $\sum b_n$ sont \textbf{absolument} convergentes alors $\sum c_n$ converge et 
\[ \sum_{n=0}^{\infty} c_n = \sum_{n=0}^{\infty} a_n \sum_{n=0}^{\infty} b_{n} . \]
\end{thm}

\begin{proof}
On suppose d'abord que les suites $(a_n)_\N$ et $(b_n)_\N$ sont positives et on note :
\[ A_N = \sum_{n=0}^{N} a_n \et B_N = \sum_{n=0}^{N} b_n \et C_N = \sum_{n=0}^{N} c_n , \]
avec $\lim A_N = A$ et $\lim B_N = B$. On note aussi les ensembles d'indices :
\[ I_n = \set{ (i,j) \in \N^2 \tq 0 \leq i, j \leq n} \et J_n = \set{ (i,j) \in \N^2 \tq 0 \leq i, j \leq n \text{ et } i+j = n} . \]
On remarque que :
\[ A_N B_N = \sum_{(i,j) \in I_N} a_i b_j  \et C_N = \sum_{(i,j) \in J_N} a_i b_j . \]
Les inclusions $I_N \subset J_{2N} \subset I_{2N}$ permettent d'affirmer que  
\[ A_N B_N \leq C_N \leq A_{2N} B_{2N} . \]
donc, d'après le théorème des gendarmes la suite $\bigpar{\sum_{k=0}^{2N} c_n}_{N \in \N}$ converge vers $AB$.
Comme de plus :
\[ \sum_{k=0}^{2N} c_n \leq \sum_{k=0}^{2N+1} c_n \leq \sum_{k=0}^{2N+2} c_n, \]
on peut bien affirmer que $\bigpar{\sum_{k=0}^{N} c_n}_{N \in \N}$ converge vers $AB$. \newline

Supposons maintenant que les suites $(a_n)_\N$ et $(b_n)_\N$ sont de signe quelconque. Alors, d'après le cas précédent, la série
de terme général $d_n = \sum_{k=0}^n \abs{a_k} \abs{b_{n-k}}$ converge. 
Or $\abs{c_n} \leq d_n$ donc la série $\sum c_n$ converge absolument.
On introduit les notations :
\[ A'_N = \sum_{n=0}^{N} \abs{a_n} \et B'_N = \sum_{n=0}^{N} \abs{b_n} \et C'_N = \sum_{n=0}^{N} \abs{c_n} , \]
on a alors :
\begin{align*}
\abs{A_N B_N - C_N} & = \abs*{\sum_{(i,j) \in I_N} a_i b_j - \sum_{(i,j) \in J_N} a_i b_j} \\
& = \abs*{\sum_{(i,j) \in I_N \setminus J_N} a_i b_j} \\
& \leq \sum_{(i,j) \in I_N \setminus J_N} \abs{a_i} \abs{b_j} \\
& \leq \sum_{(i,j) \in I_N} \abs{a_i} \abs{b_j} - \sum_{(i,j) \in J_N} \abs{a_i} \abs{b_j} \\
& \leq A'N B'N - C'N
\end{align*}
Et, en revenant au cas précédent, on a $\lim (A'N B'N - C'N) = 0$ d'où :
\[ \sum_{n=0}^{\infty} c_n = \sum_{n=0}^{\infty} a_n \sum_{n=0}^{\infty} b_{n} . \]
\end{proof}

\section{Sommation par paquets}

\begin{defn}
Soit $\sum a_n$ une série numérique et soit $\varphi : \N \to \N$ une application strictement croissante avec $\varphi(0)=0$.
Pour tout $n \in \N$, on note 
\[ b_n = \sum_{k=\varphi(n)}^{\varphi{n+1}-1} a_k . \]
 $\sum b_n$ est alors une sommation par paquets de $\sum a_n$.
\end{defn}

\begin{lem}
Si $\sum a_n$ converge alors toute sommation par paquets de $\sum a_n$ converge.
\end{lem}

\begin{thm}
Si il existe $M$ tel que $\forall n \in \N, \varphi(n+1) - \varphi(n) \leq M$ (la \og taille des paquets \fg~ est bornée) et si 
$\lim a_n = 0$ alors $\sum b_n$ converge implique que $\sum a_n$ converge.
\end{thm}

\begin{proof}
Pour tout $N \in \N$, il existe $p \in \N$ tel que $\varphi(p) \leq N < \varphi(p+1)$ donc 
\[ \abs*{\sum_{n=0}^N a_n - \sum_{n=0}^p b_n} = \abs*{\sum{n=N+1}^{\varphi(p+1)-1} a_n}, \]
avec $\sum{n=N+1}^{\varphi(p+1)-1} a_n = 0$ si $N = \varphi(p+1)-1$. Or
\[ \abs*{\sum{n=N+1}^{\varphi(p+1)-1} a_n} \leq M \cdot \sup_{\varphi(p) \leq n \leq \varphi(p+1)} (\abs{a_n}) , \]
avec $\sup_{\varphi(p) \leq n \leq \varphi(p+1)} (\abs{a_n}) \to 0$ pour $N \to \infty$.
On conclut donc que $\sum a_n \to \sum b_n$.
\end{proof}

\end{document}